\sourcepage{402}
\chapterappendix{12A}{贝叶斯临床试验}
\sourcepage{403}
多数临床试验，例如研究双膦酸盐能否降低帕金森病患者髋部骨折风险的试验，旨在给定统计学显著性水平（$\alpha$）下，检验治疗是否具有某个指定的作用，如使结局风险降低20\%。研究的样本量估计、招募与随访计划以及预算，都据此安排。研究完成后，研究者比较干预组与对照组的髋部骨折风险，计算检验统计量的$P$值，判断结果是否达到统计学显著性。如果安排中期分析，可以将$\alpha$除以中期分析次数，或以其他方式分配。这种研究设计与资料分析的\term{频率学派}（frequentist）方法，依据的是：如果目标人群中的零假设为真，样本中出现所观察到的研究结果或更极端结果的频率会有多高（见第5章）。它不纳入关于双膦酸盐对髋部骨折作用的任何既有信息。

相反，对于同一个研究问题，\term{贝叶斯方法}（Bayesian approach）（1—3）会考虑此前许多研究的结果，这些研究已表明，双膦酸盐能降低无帕金森病者的髋部骨折风险（表12A.1）。该方法首先为帕金森病患者的每一种疗效水平指定\term{先验概率}。举一个简化的例子：双膦酸盐无益甚至有害的先验概率可能为0.2，髋部骨折风险降低不超过20\%的概率为0.3，降低20\%或更多的概率为0.5；这些比例之和必须为1.0。随着研究中髋部骨折的发生，将观察资料与这些先验概率结合，便得到新的可能性，称为\term{后验概率}（posterior probabilities）。临床试验的贝叶斯方法，类似于临床医师将疾病的检验前概率与诊断检验结果结合，以确定疾病检验后概率的做法（4）。

\sourcepage{404}
\begin{longtable}{L{27mm}L{61mm}L{61mm}}
\caption*{\textbf{表12A.1\quad 临床试验的频率学派与贝叶斯方法}}\\\toprule
概念&频率学派方法&贝叶斯方法\\\midrule\endfirsthead
\toprule 概念&频率学派方法&贝叶斯方法\\\midrule\endhead
试验结果&用点估计和置信区间表示治疗作用&治疗作用的后验分布、后验均数或中位数，以及可信区间\\
对既有研究资料的考虑&即使使用，也只是主观地用于解释$P$值&分析必须指定疗效的先验概率分布；完全没有既有资料时，通常采用平坦先验或怀疑先验（见正文）\\
统计学显著性或充分疗效依据&零假设的$P$值小于预先规定的阈值，通常为0.05&疗效的后验概率较高；申请新治疗的监管批准时，要求可能更严格\\
关于疗效的临床专家意见的作用&无&用于设定先验概率分布，通常涉及判断既有研究结果在多大程度上适用于当前试验\\
主观决定&将$P$值解释为存在治疗作用的可能性；规定样本量估计所需的参数；规定资料监察方式；选择分析模型&选择先验概率分布；选择分析模型\\
样本量&根据效应量、$\alpha$和$\beta$预先规定；效应量通常是尽可能合理的猜测&可以不预先规定；持续纳入参与者，直至有效、有害或无效的概率达到某一水平。有时也通过估计获得充分资料所需的预定样本量来设定\\
预算&固定样本量使预算可以固定&完全序贯的贝叶斯设计需要灵活的预算\\
资料监察与停止规则&在指定时点分析资料；基于$P$值的停止规则要对中期分析次数作调整&结果可不断更新，无须因反复分析而作惩罚性调整；新的后验概率取代原来的后验概率\\\bottomrule
\end{longtable}
\sourcepage{405}
治疗有效的先验概率不是一个数，甚至也不是前例中的几个数，而是表示为\term{先验分布}（prior distribution）。它通常是一条正态分布的钟形曲线，将疗效作为连续变量，估计不同疗效出现的可能性（图12A.1）。当我们对治疗作用知道得很多时，先验分布较窄；既有资料有限时，分布较宽。理想情况下，先验分布应依据此前针对类似研究问题的高质量研究结果；但有时必须依赖专家意见与主观判断。
\sourcepage{406}
\begin{figure}[htbp]\centering\chTwelvePriors
\caption*{\textbf{图12A.1\quad 先验分布示例。}横轴表示治疗的疗效，定义为$1-\mathrm{RR}$，其中RR为治疗相对于对照的风险比；纵轴表示该疗效的先验概率。红线为较窄的先验分布，以15\%的疗效为中心；绿线为较宽的先验分布，同样以15\%为中心；蓝线为以0\%疗效为中心的怀疑先验；紫线为平坦的无信息先验，也以0\%为中心。每条曲线下的面积均为1.0。}
\end{figure}

如何确定先验分布呢？例如，一项新治疗只是对既有药物略作修改，而数项实施良好的试验已发现，该既有药物可将所关注结局的风险降低约15\%。此时，新治疗很可能有相近效果，可以据此采用位于有效一侧的乐观先验分布，例如获益范围为7.5\%～22.5\%。另一方面，如果此前类似药物的研究有限，例如只研究过对生物标志物的作用，那么\sourcepage{407}采用以零为中心、认为治疗带来巨大益处或危害的可能性都很小的怀疑先验分布，就比较合适。先验分布通常呈钟形，但无信息先验可以是平坦的，即在所有合理的获益或危害水平上均匀分布，以表示缺乏信息。

从既有研究结果得到的先验分布，常常还要与怀疑先验或无信息先验混合，从而加以修正。混合比例取决于专家判断：既有研究适用于当前临床情境的概率有多大。这样便生成一个新的先验分布。

分析资料时，贝叶斯试验不检验试验结果是否具有统计学显著性，而是将结果与先验分布结合，生成\term{后验分布}（posterior distribution），表示治疗作用可能落入的范围；通常用\term{可信区间}（credible interval）或\term{最高后验密度区间}（highest posterior density interval）概括。例如，假定所选先验分布有效，治疗作用落在90\%可信区间内的概率就是0.9。也可以将结果表示为某个治疗作用的概率，例如治疗使结局风险至少降低30\%的概率为0.94；这一表述同样以所选先验分布为前提。

贝叶斯方法可能得出与频率学派不同的结论。例如，Brophy（5）重新分析了一项试验，该试验将左主干冠状动脉疾病患者随机分配到经皮冠状动脉介入治疗（percutaneous coronary intervention，PCI）或冠状动脉旁路移植术（coronary artery bypass graft，CABG）。研究发现，PCI患者的主要结局，即5年内死亡、心肌梗死和卒中，明显多于CABG患者，但风险增加的95\%置信区间（$-0.9\%$～6.5\%）包含零。因此，采用频率学派方法时，研究结论是PCI与CABG的结局没有显著差异，这一结论能给读者传达的信息很有限。相反，贝叶斯方法纳入此前研究类似问题的试验结果后，估计PCI增加试验终点风险的概率为0.96。

\sourcepage{408}
贝叶斯试验的设计也在若干方面不同于经典随机试验。有时，可以估计获得充分资料、形成“足够窄”的后验分布所需的预期参与者人数，以此估计样本量，但不必预先固定。相反，可以采用序贯方法，随着结局资料不断收集，更新治疗有效或有害的概率。由于贝叶斯方法不依赖$P$值与$\alpha$，因此无须对多次查看资料作调整（6）。当该概率超过预先设定的有效、有害或徒劳阈值时，可以停止招募；徒劳意味着研究结果不大可能回答研究问题。当然，贝叶斯试验的研究者必须指定疗效的先验分布，规定如何决定是否修改该分布，例如是否与怀疑先验混合，还要规定研究期间是否更新这些决定，以及如何更新。

没有固定样本量的设计，有时称为“完全序贯”设计。由于研究开始时可能尚不知道样本量和随访时间，因此需要灵活的资助安排。资助者不能只拨出一笔固定预算，而必须接受研究可能比希望的时间更长、费用更多；当然，研究也可能比预期更早结束。

总的来说，贝叶斯试验在设计和分析时纳入关于干预疗效的既有资料。当既往资料足以帮助确定先验分布时，贝叶斯方法可能合适。如果观察结果与先验预期一致，贝叶斯试验有可能需要更小的样本量。此外，将研究结果表示为治疗疗效的一系列概率，可能比单个$P$值或置信区间提供更多信息。不过，贝叶斯方法依赖对先验分布的主观判断，其中包括既有资料在多大程度上适用于目标人群和所研究的治疗。它还要求预算具有灵活性，而美国国立卫生研究院（National Institutes of Health，NIH）等传统资助来源可能无法提供这种安排。
\sourcepage{409}
\begin{references}
\item Harrell FE, Jr. Introduction to Bayes for Evaluating Treatments. http://hbiostat.org/doc/bayes/course.html
\item Diamond GA, Kaul S. Prior convictions: Bayesian approaches to the analysis and interpretation of clinical megatrials. J Am Coll Cardiol. 2004;43(11):1929-1939.
\item Berry DA. Bayesian clinical trials. Nat Rev Drug Discov. 2006;5:27-36.
\item Browner WS, Newman TB. Are all significant P values created equal? The analogy between diagnostic tests and clinical research. JAMA. 1987;257:2459-2463.
\item Brophy JM. Bayesian interpretation of the EXCEL trial and other randomized clinical trials of left main coronary artery revascularization. JAMA Intern Med. 2020;180:986-992.
\item Gelman A, Hill J, Yajima M. Why we (usually) don’t have to worry about multiple comparisons. J Res Educ Effective. 2012;5:189-211.
\end{references}


\chapterappendix{12B}{第12章习题：其他干预性研究设计}
\sourcepage{410}
\begin{enumerate}
\item 非那雄胺（finasteride）是一种5α还原酶抑制剂，也是少数获得美国食品药品监督管理局（Food and Drug Administration，FDA）批准、用于治疗男性型脱发的药物之一。每日使用时，它对促进头发再生有中等程度的疗效，但停药6个月后作用即消失。另一类药物——他汀类——已被发现能通过与非那雄胺不同的生物学通路，增加啮齿类动物的毛发生长。假设一家初创公司想收集资料，证明一种新型外用他汀类药物HairStat治疗男性型脱发的疗效。针对以下每一种试验方法，至少列出一个优点和一个缺点。
\begin{enumerate}
\item 将男性型脱发者随机分配到HairStat或安慰剂组，两组均不提供非那雄胺。
\item 开展HairStat与非那雄胺的随机试验，比较两种治疗的疗效。
\item 将男性随机分配到HairStat或安慰剂组，但只招募此前使用非那雄胺未见获益，或不能耐受非那雄胺的人。
\item 采用析因试验设计，将男性随机分配到四组：（1）非那雄胺与HairStat联合；（2）非那雄胺与HairStat安慰剂；（3）非那雄胺安慰剂与HairStat；（4）双安慰剂。
\item 开展交叉试验，将男性随机分配到先用非那雄胺、3个月后改用HairStat，或按相反顺序用药的两组。
\end{enumerate}
\item 你在一家设有多个办公地点的大公司工作，公司希望减少员工肥胖，并促进健康饮食习惯。一位顾问建议实施多成分饮食干预，包括发放营养宣传册、在工作场所食堂开展营养宣传、在公共区域张贴教育海报，以及按照轮班时间安排饮食意识工作坊。你想评估这项干预能否促进超重或肥胖员工减重。以下实施和评估方法各有哪些潜在优点或缺点？
\begin{enumerate}
\item 在全公司范围内，将员工个人随机分配到饮食支持干预或无干预组，并评估6个月内每个人的体重变化。
\item 在公司的所有办公地点同时实施干预，持续6个月，比较员工干预前后的体重变化。\sourcepage{411}
\item 开展阶梯楔形试验，在6个月内按随机顺序，依次将饮食干预引入公司的所有办公地点。
\item 开展整群随机试验，将公司的各办公地点随机分配到实施饮食支持干预或不实施干预的两组。
\end{enumerate}
\end{enumerate}
